%!TEX root = ../../main.tex
% Section 3: Method

\section{Method}\label{sec:method}

We present the Autoencoding Factor Model (AFM), a neural architecture that learns tradable factor portfolios from mutual fund returns and fund-level characteristics. The key innovation is a \emph{portfolio bottleneck}: each latent factor is defined as the return of an explicit portfolio whose weights are produced by a neural network using only past information, guaranteeing tradability by construction. We formalize the problem setup (\S\ref{sec:setup}), describe the architecture (\S\ref{sec:architecture}--\S\ref{sec:subnetworks}), present the training objective (\S\ref{sec:training}), and discuss the role of the portfolio bottleneck (\S\ref{sec:discussion}).

\subsection{Problem Setup}\label{sec:setup}

We observe a panel of excess returns for $N$ mutual funds over $T$ trading days. Let $r_{i,t} \in \mathbb{R}$ denote the excess return of fund $i$ at time $t$, and stack these as $\mathbf{r}_t = (r_{1,t}, \ldots, r_{N,t})^\top \in \mathbb{R}^N$. Each fund $i$ has an observable characteristic vector $\mathbf{x}_{i,t} \in \mathbb{R}^D$ derived from regulatory filings (economic sector holdings and asset class allocations), with the cross-section stacked as $\mathbf{X}_t \in \mathbb{R}^{N \times D}$. Fix a lookback window of length $L$; the lagged return panel available at time $t-1$ is
\begin{equation}\label{eq:lookback}
    \mathbf{R}_{t-1}^{(L-1)} := \big(\mathbf{r}_{t-L+1}, \ldots, \mathbf{r}_{t-1}\big) \in \mathbb{R}^{N \times (L-1)}.
\end{equation}
The information set at time $t-1$ is $\mathcal{F}_{t-1} := \sigma\!\left(\mathbf{X}_{t-1}, \mathbf{R}_{t-1}^{(L-1)}\right)$. All quantities used to form portfolios at time $t$ are functions of $\mathcal{F}_{t-1}$, ensuring \textbf{no look-ahead}.

\paragraph{Goal.} Learn $F$ latent factors that admit a conditional linear factor model
\begin{equation}\label{eq:factor-model}
    \mathbf{r}_t \approx \boldsymbol{\alpha}_{t-1} + \mathbf{B}_{t-1} \mathbf{f}_t,
\end{equation}
where $\boldsymbol{\alpha}_{t-1} \in \mathbb{R}^N$ is a vector of asset intercepts, $\mathbf{B}_{t-1} \in \mathbb{R}^{N \times F}$ is a matrix of conditional factor loadings, and $\mathbf{f}_t \in \mathbb{R}^F$ is a vector of realized factor returns---all with $\boldsymbol{\alpha}_{t-1}$ and $\mathbf{B}_{t-1}$ measurable with respect to $\mathcal{F}_{t-1}$. Crucially, we require that factors are \emph{tradable portfolios}:
\begin{equation}\label{eq:tradable}
    \mathbf{f}_t := \mathbf{W}_{t-1}^\top \mathbf{r}_t, \qquad \mathbf{W}_{t-1} \in \mathbb{R}^{N \times F},
\end{equation}
where $\mathbf{W}_{t-1}$ is $\mathcal{F}_{t-1}$-measurable. Table~\ref{tab:notation} summarizes the notation used throughout.

\begin{table}[t]
\centering
\caption{Notation summary.}
\label{tab:notation}
\small
\begin{tabular}{cl}
\toprule
\textbf{Symbol} & \textbf{Meaning} \\
\midrule
$N, T, F$ & Number of funds, time periods, latent factors \\
$D, S, L$ & Characteristic dim., state dim., lookback length \\
$r_{i,t}$, $\mathbf{r}_t$ & Excess return of fund $i$; $N$-vector of returns \\
$\mathbf{x}_{i,t}$, $\mathbf{X}_t$ & Characteristic vector of fund $i$; $N \times D$ matrix \\
$\mathbf{R}_{t-1}^{(L-1)}$ & $N \times (L-1)$ lagged return panel \\
$\mathcal{F}_{t-1}$ & Information set at $t-1$ \\
$\mathbf{s}_{t-1}$ & Learned market state vector ($\in \mathbb{R}^S$) \\
$\mathbf{W}_{t-1}$, $\mathbf{w}_{t-1}^{(k)}$ & $N \times F$ weight matrix; weights for factor $k$ \\
$\mathbf{f}_t$ & $F$-vector of realized factor returns \\
$\mathbf{B}_{t-1}$, $\boldsymbol{\beta}_{i,t-1}$ & $N \times F$ loading matrix; loadings for fund $i$ \\
$\boldsymbol{\alpha}_{t-1}$ & $N$-vector of asset intercepts \\
$g_\psi, h_\phi, b_\theta, a_\eta$ & ReturnEncoder, ScoreNet, BetaNet, AlphaNet \\
$k_0$ & Number of long-only factors \\
\bottomrule
\end{tabular}
\end{table}


\subsection{Architecture Overview}\label{sec:architecture}

The AFM has the structure of an autoencoder whose bottleneck is a vector of portfolio returns (Figure~\ref{fig:architecture}). The \emph{encoder} maps the information set $\mathcal{F}_{t-1}$ to portfolio weights $\mathbf{W}_{t-1}$ and factor exposures $(\boldsymbol{\alpha}_{t-1}, \mathbf{B}_{t-1})$. The \emph{bottleneck} computes factor returns $\mathbf{f}_t = \mathbf{W}_{t-1}^\top \mathbf{r}_t$ as realized portfolio returns. The \emph{decoder} reconstructs fund returns via $\hat{\mathbf{r}}_t = \boldsymbol{\alpha}_{t-1} + \mathbf{B}_{t-1} \mathbf{f}_t$.

\begin{figure}[t]
\centering
\fbox{\rule{0pt}{2in}\rule{0.9\linewidth}{0pt}} % [PLACEHOLDER] Replace with figures/architecture.pdf
\caption{\textbf{AFM architecture overview.} The encoder maps fund characteristics and a learned market-state embedding to portfolio weights (ScoreNet), factor loadings (BetaNet), and intercepts (AlphaNet). The bottleneck computes factor returns as explicit portfolio returns $\mathbf{f}_t = \mathbf{W}_{t-1}^\top \mathbf{r}_t$. The decoder reconstructs fund returns via $\hat{\mathbf{r}}_t = \boldsymbol{\alpha}_{t-1} + \mathbf{B}_{t-1} \mathbf{f}_t$. All weights and loadings depend only on $\mathcal{F}_{t-1}$, ensuring no look-ahead.}
\label{fig:architecture}
\end{figure}

Four sub-networks parameterize this mapping:
\begin{enumerate}[leftmargin=*,nosep]
    \item \textbf{ReturnEncoder} ($g_\psi$): compresses lagged returns into a market state vector;
    \item \textbf{ScoreNet} ($h_\phi$): produces portfolio scores, mapped to valid weights via softmax;
    \item \textbf{BetaNet} ($b_\theta$): outputs conditional factor loadings per fund;
    \item \textbf{AlphaNet} ($a_\eta$): outputs fund-specific intercepts.
\end{enumerate}
All four networks receive only $\mathcal{F}_{t-1}$-measurable inputs, so every output---weights, loadings, and intercepts---satisfies the no-look-ahead requirement by construction. We describe each sub-network next.


\subsection{Sub-Network Details}\label{sec:subnetworks}

\paragraph{ReturnEncoder.} Financial cross-sections are shaped by time-varying market conditions---volatility regimes, interest rate environments, and risk appetite. The ReturnEncoder compresses the lagged return panel into a low-dimensional state vector:
\begin{equation}\label{eq:return-encoder}
    \mathbf{s}_{t-1} = g_\psi\!\left(\operatorname{vec}\!\left(\mathbf{R}_{t-1}^{(L-1)}\right)\right) \in \mathbb{R}^{S},
\end{equation}
where $g_\psi$ is an MLP with layer normalization and GELU activations. The state vector $\mathbf{s}_{t-1}$ is broadcast across all funds when computing per-fund outputs, enabling the same characteristic vector $\mathbf{x}_{i,t-1}$ to produce different portfolio weights and exposures under different market conditions without hand-crafted regime indicators.

\paragraph{ScoreNet and portfolio weight construction.} The ScoreNet produces a raw score for each fund-factor pair:
\begin{equation}\label{eq:scorenet}
    \operatorname{score}_{i,t-1}^{(k)} = h_\phi\!\left([\mathbf{x}_{i,t-1};\, \mathbf{s}_{t-1}]\right)_k, \qquad k = 1, \ldots, F,
\end{equation}
where $[\,\cdot\,;\,\cdot\,]$ denotes concatenation. We convert scores to valid portfolio weights using two differentiable mappings, depending on the factor type.

\emph{Long-only factors} ($k = 1, \ldots, k_0$). Scores are passed through a cross-sectional softmax:
\begin{equation}\label{eq:long-only}
    w_{i,t-1}^{(k)} = \frac{\exp\!\left(\operatorname{score}_{i,t-1}^{(k)}\right)}{\sum_{j=1}^{N} \exp\!\left(\operatorname{score}_{j,t-1}^{(k)}\right)},
\end{equation}
guaranteeing $w_{i,t-1}^{(k)} \geq 0$ and $\sum_{i} w_{i,t-1}^{(k)} = 1$ by construction. The resulting factor $f_t^{(k)}$ is the return of a fully-invested long-only portfolio.

\emph{Long-short market-neutral factors} ($k = k_0+1, \ldots, F$). We employ a softmax-minus-softmax construction:
\begin{equation}\label{eq:long-short}
    w_{i,t-1}^{(k)} = \frac{\exp\!\left(\operatorname{score}_{i,t-1}^{(k)}\right)}{\sum_{j} \exp\!\left(\operatorname{score}_{j,t-1}^{(k)}\right)} - \frac{\exp\!\left(-\operatorname{score}_{i,t-1}^{(k)}\right)}{\sum_{j} \exp\!\left(-\operatorname{score}_{j,t-1}^{(k)}\right)}.
\end{equation}
Since each softmax sums to one, market neutrality holds exactly: $\sum_{i=1}^{N} w_{i,t-1}^{(k)} = 1 - 1 = 0$. We further normalize by gross leverage, $\mathbf{w}_{t-1}^{(k)} \leftarrow \mathbf{w}_{t-1}^{(k)} \big/ \sum_{i} |w_{i,t-1}^{(k)}|$, to stabilize factor scaling across training epochs. This construction is a differentiable analogue of ``long the highest-score funds, short the lowest-score funds,'' avoiding non-differentiable sorting operations while enforcing self-financing structure.

The final weight matrix concatenates long-only and long-short columns:
\begin{equation}\label{eq:weight-matrix}
    \mathbf{W}_{t-1} = \big[\mathbf{W}_{t-1}^{\mathrm{LO}} \;\; \mathbf{W}_{t-1}^{\mathrm{LS}}\big] \in \mathbb{R}^{N \times F}.
\end{equation}
In our experiments, we use $k_0 = 2$ long-only factors and $F - k_0 = 3$ long-short factors.

\paragraph{BetaNet.} For each fund $i$, the BetaNet outputs conditional factor loadings:
\begin{equation}\label{eq:betanet}
    \boldsymbol{\beta}_{i,t-1} = b_\theta\!\left([\mathbf{x}_{i,t-1};\, \mathbf{s}_{t-1}]\right) \in \mathbb{R}^{F},
\end{equation}
stacked row-wise into the loading matrix $\mathbf{B}_{t-1} \in \mathbb{R}^{N \times F}$. These loadings are the model's analogue of classical factor betas: $\beta_{i,t-1}^{(k)}$ measures fund $i$'s exposure to factor $k$, conditional on current holdings and market state, enabling risk attribution and portfolio construction.

\paragraph{AlphaNet.} The AlphaNet produces fund-specific intercepts:
\begin{equation}\label{eq:alphanet}
    \alpha_{i,t-1} = a_\eta\!\left([\mathbf{x}_{i,t-1};\, \mathbf{s}_{t-1}]\right) \in \mathbb{R},
\end{equation}
capturing residual expected return not explained by factor exposures.

\paragraph{Reconstruction.} Given weights $\mathbf{W}_{t-1}$ and realized returns $\mathbf{r}_t$, factor returns and reconstructed returns are computed as:
\begin{equation}\label{eq:reconstruction}
    \mathbf{f}_t = \mathbf{W}_{t-1}^\top \mathbf{r}_t \in \mathbb{R}^F, \qquad \hat{\mathbf{r}}_t = \boldsymbol{\alpha}_{t-1} + \mathbf{B}_{t-1} \mathbf{f}_t \in \mathbb{R}^N.
\end{equation}
The bottleneck $\mathbf{f}_t$ has dimension $F \ll N$; training encourages portfolio weights that produce factor returns maximally informative for reconstructing the entire cross-section through the linear decoder $\mathbf{B}_{t-1}$.


\subsection{Training Objective}\label{sec:training}

Let $\Theta = (\psi, \phi, \theta, \eta)$ denote all learnable parameters. We minimize the composite objective:
\begin{equation}\label{eq:total-loss}
    \mathcal{L}(\Theta) = \mathbb{E}_{t}\!\left[\ell_{\mathrm{rec}}(t) + \lambda_{\mathrm{orth}}\,\ell_{\mathrm{orth}}(t) + \lambda_{\beta}\,\ell_{\beta}(t) + \lambda_{w}\,\ell_{w}(t)\right],
\end{equation}
where expectations are approximated by mini-batch averages over time indices.

\paragraph{Reconstruction loss.} The primary objective is cross-sectional mean squared error:
\begin{equation}\label{eq:rec-loss}
    \ell_{\mathrm{rec}}(t) = \frac{1}{N}\sum_{i=1}^{N} \left(\hat{r}_{i,t} - r_{i,t}\right)^2.
\end{equation}

\paragraph{Orthogonality penalty.} Without explicit encouragement, multiple factor portfolios may collapse onto the same dominant direction (e.g., a market-like portfolio). Let $\widehat{\boldsymbol{\Sigma}}_f$ denote the sample covariance matrix of factor returns $\{\mathbf{f}_t\}$ within a mini-batch. We penalize off-diagonal mass:
\begin{equation}\label{eq:orth-loss}
    \ell_{\mathrm{orth}} = \left\|\widehat{\boldsymbol{\Sigma}}_f - \operatorname{diag}\!\left(\widehat{\boldsymbol{\Sigma}}_f\right)\right\|_F^2,
\end{equation}
encouraging the model to discover distinct sources of variation.

\paragraph{Regularization.} We add $\ell_2$ penalties on factor loadings and portfolio weights to control model capacity and prevent extreme concentrations:
\begin{equation}\label{eq:reg}
    \ell_{\beta}(t) = \|\mathbf{B}_{t-1}\|_F^2, \qquad \ell_{w}(t) = \|\mathbf{W}_{t-1}\|_F^2.
\end{equation}
These penalties shrink extreme exposures and overly concentrated portfolio weights, improving out-of-sample robustness.

\paragraph{Optimization.} We train with Adam \citep{Kingma2015Adam} using a learning rate of $10^{-2}$ and batch size 128 for 150 epochs. The best model state is selected by validation loss (computed on the last 10\% of dates) with early stopping. Full hyperparameter details are in Appendix~\ref{app:implementation}.


\subsection{Discussion: Why the Portfolio Bottleneck Matters}\label{sec:discussion}

The portfolio bottleneck is what distinguishes the AFM from standard conditional autoencoders for asset pricing. In the model of \citet{GuKellyXiu2021autoencoder}, latent factors $\mathbf{f}_t$ are free variables produced by the encoder---they have no direct portfolio interpretation and cannot be traded without an additional projection step. By contrast, our constraint $\mathbf{f}_t = \mathbf{W}_{t-1}^\top \mathbf{r}_t$ provides three advantages.

\paragraph{Tradability.} Each factor $f_t^{(k)} = {\mathbf{w}_{t-1}^{(k)}}^\top \mathbf{r}_t$ is the realized return of a portfolio with known, pre-determined weights. A practitioner can implement these factors as live trading strategies---long-only or market-neutral---without any post-hoc approximation.

\paragraph{Interpretability.} The portfolio weights $\mathbf{w}_{t-1}^{(k)}$ reveal which funds drive each factor: inspecting the largest weights identifies the economic themes each factor captures. The conditional loadings $\boldsymbol{\beta}_{i,t-1}$ provide complementary fund-level exposure diagnostics for risk attribution.

\paragraph{Implicit regularization.} Restricting factors to the span of realizable portfolio returns constrains the model's hypothesis class relative to free latent codes. The model cannot exploit arbitrary rotations of the latent space to minimize reconstruction error; it must find directions in \emph{portfolio-return space} that are genuinely informative for the cross-section. This connection echoes the classical theory of mimicking portfolios \citep{Huberman1987mimicking}, where any priced factor admits a portfolio representation---the AFM learns such portfolios directly and nonlinearly. The portfolio bottleneck also relates to differentiable portfolio optimization \citep{Donti2017taskbased, Agrawal2019differentiable}, but targets cross-sectional explanatory power rather than portfolio performance. The softmax weight construction can be viewed as a differentiable relaxation of parametric portfolio policies \citep{Brandt2009parametric}, extended from a single-portfolio objective to a multi-factor autoencoding framework.
